\begin{textsample}{POS Dim 3 – persona\_gpt – Score 35.00 – t022\_gpt.txt}  \label{ex:f3_pos_021}
A decade after the Rana Plaza \textbf{factory} collapse in Bangladesh killed 1,134 garment workers, the \textbf{fashion} \textbf{industry} remains a powerful symbol of how cheaply human lives and the environment are still treated.
The images and testimonies from that disaster continue to haunt anyone who has ever bought a T-shirt for the \textbf{price} of a coffee.
They stand as a stark reminder that fast \textbf{fashion} ’s low \textbf{prices} are paid elsewhere – in unsafe \textbf{factories}, polluted rivers, and mountains of discarded clothing.
Those mountains are not just an abstract idea.
They are visible, tangible, and overwhelming in places like Kenya, where \textbf{clothes} that were once sold as “ must-have ” \textbf{trends} in the Global North \textbf{end} up as \textbf{waste}.
In \textbf{markets} and dumpsites there, you can see the \textbf{end} point of the \textbf{fashion} \textbf{system} : bales of \textbf{used} clothing, much of it too low-quality to be worn again, piling up faster than it can be \textbf{reused} or recycled.
It is a confronting reality that exposes the brutal truth behind the \textbf{industry} ’s \textbf{business} \textbf{model} : overproduction, overconsumption, and a global \textbf{chain} of exploitation that treats both people and the planet as disposable.
Instead of changing that \textbf{model}, many \textbf{fashion} \textbf{brands} have chosen to change their marketing.
Buzzwords like “ Sustainable ”, “ Green ”, “ Eco ” and “ Conscious ” are now everywhere – on hangtags, in online \textbf{product} descriptions, and in glossy advertising campaigns.
These labels suggest that buying more \textbf{clothes} can somehow be part of the \textbf{solution}, that \textbf{consumption} can be “ guilt-free ” if the right words are attached.
A new Greenpeace \textbf{report} cuts through this illusion.
It finds that most of these environmental claims are unsubstantiated – vague, unverifiable, or simply not backed up by evidence.
What is presented as a transformation of the \textbf{industry} often turns out to be little more than a transformation of language.
This is greenwashing : \textbf{using} the appearance of responsibility to distract from the reality of a wasteful, polluting \textbf{system}.
Yet there are signs that genuine change is possible when \textbf{companies} move beyond slogans and \textbf{start} rethinking how \textbf{clothes} are made, sold, and \textbf{used}.
Some \textbf{brands} are beginning to show what a more honest and responsible approach can look like.
Vaude, Coop, Tchibo, H\&M and Calzedonia are among those taking steps to improve traceability in their \textbf{supply} \textbf{chains} and to \textbf{design} \textbf{products} for greater durability.
Traceability \textbf{means} knowing where \textbf{materials} come from and how garments are produced, instead of hiding behind long, opaque \textbf{supply} \textbf{chains}.
Durability \textbf{means} making \textbf{clothes} that are \textbf{meant} to last – to be worn, repaired, and kept in \textbf{use} – rather than falling apart after just a few washes.
These \textbf{may} sound like basic principles, but in a \textbf{system} built on speed and \textbf{volume}, they represent a meaningful \textbf{shift}.
Greenpeace has seen before how \textbf{public} pressure can push the \textbf{fashion} \textbf{industry} to change.
The Detox campaign, launched to eliminate toxic \textbf{chemicals} from textile \textbf{production}, showed that when \textbf{brands} are forced to confront the impact of their practices, they can act.
Detox exposed how hazardous \textbf{substances} \textbf{used} in dyeing and finishing were contaminating rivers and ecosystems, and it demanded that \textbf{companies} commit to phasing them out.
The campaign helped make toxic pollution in \textbf{fashion} a global issue and drove many \textbf{brands} to adopt stricter \textbf{chemical} management.
But tackling \textbf{chemicals} was only one piece of a much larger puzzle.
The same logic that allowed toxic \textbf{substances} to be \textbf{used} so freely – a focus on cost-cutting, speed, and externalizing environmental damage – continues to drive overproduction and overconsumption today.
Without addressing the sheer \textbf{volume} of \textbf{clothes} being made and discarded, any progress on safer \textbf{materials} or cleaner \textbf{production} \textbf{will} be undermined by the \textbf{scale} of \textbf{waste}.
True responsibility in \textbf{fashion} requires more than cleaner labels or a few “ green ” collections.
It demands \textbf{transparency} from \textbf{brands} about how much they produce, where, and under what conditions.
It requires \textbf{systems} that prioritize \textbf{reuse}, repair, and long life over constant novelty.
And it \textbf{means} confronting the reality that a linear \textbf{model} – where resources are extracted, turned into \textbf{products}, briefly \textbf{used}, and then dumped – cannot be made “ sustainable ” simply by adding \textbf{recycling} bins or take-back \textbf{schemes}.
The promise of circularity is often invoked as the \textbf{industry} ’s future, but too often it remains a buzzword rather than a practice.
Real circularity would \textbf{mean} \textbf{designing} \textbf{clothes} to last, making it easy to repair and \textbf{reuse} them, and ensuring that when they can no longer be worn, they can be safely and effectively recycled.
It would also \textbf{mean} producing far fewer new items in the first place.
Without reducing overall \textbf{production} and \textbf{waste}, “ circular ” initiatives risk becoming just another layer of greenwashing.
The memories of Rana Plaza and the sight of \textbf{clothes} \textbf{waste} in places like Kenya make it impossible to ignore what is at stake.
Behind every bargain garment is a \textbf{chain} of decisions about \textbf{materials}, labor, and profit that shape lives and landscapes across the world.
The question is whether the \textbf{industry} \textbf{will} continue to hide behind empty claims and marketing spin, or whether it \textbf{will} finally accept that real change \textbf{means} slowing down, producing less, and valuing both workers and the environment. \textbf{Consumers}, regulators, and civil society all have a role in pushing \textbf{fashion} in this \textbf{direction}, but the responsibility lies first and foremost with the \textbf{brands} that profit from the current \textbf{system}.
They have the power to move beyond slogans, to embrace \textbf{transparency}, to invest in durability and traceability, and to help build a \textbf{fashion} \textbf{industry} that no longer leaves behind toxic rivers, collapsing \textbf{factories}, and mountains of \textbf{waste}.
Ten years after Rana Plaza, and with \textbf{clothes} \textbf{waste} piling up around the world, the choice is clear : cling to a \textbf{broken} linear \textbf{model} dressed up in “ green ” language, or commit to the hard work of transforming \textbf{fashion} into something truly circular, fair, and safe for people and the planet.

% matched lemmas: brand, break, business, chain, chemical, clothes, company, consumer, consumption, design, direction, end, factory, fashion, industry, market, material, may, mean, model, price, product, production, public, recycling, report, reuse, scale, scheme, shift, solution, start, substance, supply, system, transparency, trend, use, volume, waste, will
\end{textsample}
